\documentclass[11pt,a4paper]{article}
\usepackage[margin=20mm,headheight=15pt]{geometry}
\usepackage{fontspec}
\setmainfont{TeX Gyre Pagella}
\setmonofont{DejaVu Sans Mono}[Scale=0.82]
\usepackage{amsmath}
\usepackage{unicode-math}
\setmathfont{TeX Gyre Pagella Math}
\setcounter{MaxMatrixCols}{20}
\usepackage{microtype}
\usepackage{enumitem}
\usepackage{fvextra}
\usepackage{booktabs}
\usepackage{longtable}
\usepackage{fancyhdr}
\usepackage[hidelinks]{hyperref}
\setlength{\parindent}{0pt}
\setlength{\parskip}{6pt}
\setlength{\emergencystretch}{3em}
\hyphenpenalty=10000
\exhyphenpenalty=10000
\raggedright
\DefineVerbatimEnvironment{Code}{Verbatim}{fontsize=\small,
  breaklines=true,breakanywhere=true,breaksymbolleft={},breaksymbolright={},
  frame=single,framesep=5pt}
\setlist[enumerate]{leftmargin=7mm,itemsep=3pt,topsep=5pt,parsep=0pt}
\pagestyle{fancy}
\fancyhf{}
\fancyhead[L]{\small Computer Vision 1}
\fancyhead[R]{\small Exam 1 · Solutions}
\fancyfoot[C]{\small \thepage}
\hypersetup{pdftitle={Exam 1 — Solutions},pdfsubject={CV1 practice material}}
\begin{document}
{\LARGE\bfseries Exam 1 — Solutions}\par
{\large Foundations under pressure}\par\bigskip
Answer key and worked explanations. Source aliases and convention warnings are listed in \texttt{COVERAGE.md} and \texttt{SOURCE\_NOTES.md} in the study pack.
\section*{Quick key}
\begin{longtable}{@{}r l@{\hspace{25mm}}r l@{}}
\toprule Question & Answer & Question & Answer \\ \midrule
01 & B & 26 & A \\
02 & A & 27 & D \\
03 & A & 28 & C \\
04 & A & 29 & A \\
05 & C & 30 & B \\
06 & D & 31 & -270.00 \\
07 & C & 32 & 0.75 \\
08 & B & 33 & 4.00 \\
09 & C & 34 & 1 \\
10 & A & 35 & 8 \\
11 & C & 36 & 8 \\
12 & B & 37 & 2.83 \\
13 & A & 38 & 59.04 \\
14 & D & 39 & 2.36 \\
15 & B & 40 & 34 \\
16 & D & 41 & 7.33 \\
17 & C & 42 & 48.75 \\
18 & A & 43 & 0.63 \\
19 & D & 44 & 0.57 \\
20 & D & 45 & 0.33 \\
21 & C & 46 & 1680 \\
22 & B & 47 & 368 \\
23 & B & 48 & 7850 \\
24 & B & 49 & 8 \\
25 & D & 50 & 2.00 \\
\bottomrule\end{longtable}
\clearpage\section*{Worked solutions}
\begin{minipage}{\linewidth}
\textbf{Q01 — B}\par\smallskip
\textbf{B. \(C=-R^{\mathsf T}t\)}\par
At the optical centre \(X_c=0\). Thus \(R C+t=0\) and \(C=-R^{-1}t=-R^{\mathsf T}t\). The translation column is not generally the camera position.\par
{\small\itshape Topic: camera. Sources: HC1b pp.35–45; S1.}

\end{minipage}
\par\vspace{12pt}
\begin{minipage}{\linewidth}
\textbf{Q02 — A}\par\smallskip
\textbf{A. Halve \(f_x\),\allowbreak{} \(f_y\),\allowbreak{} skew,\allowbreak{} \(c_x\) and \(c_y\); keep R and t unchanged.}\par
Pixel coordinates change,\allowbreak{} not the camera pose. The first two rows of K scale by 1/\allowbreak{}2 and its last row stays \(\begin{bmatrix}0&0&1\end{bmatrix}\). Scaling the entire K only changes homogeneous scale.\par
{\small\itshape Topic: camera. Sources: HC1b pp.35–45; S1.}

\end{minipage}
\par\vspace{12pt}
\begin{minipage}{\linewidth}
\textbf{Q03 — A}\par\smallskip
\textbf{A. They lie on the same viewing ray from the optical centre.}\par
Projection retains the direction of a camera-centred ray but loses distance along it. This is a fundamental single-view depth ambiguity.\par
{\small\itshape Topic: camera. Sources: HC1a; HC1b pp.24–34.}

\end{minipage}
\par\vspace{12pt}
\begin{minipage}{\linewidth}
\textbf{Q04 — A}\par\smallskip
\textbf{A. Its outgoing radiance,\allowbreak{} apart from visibility and camera effects.}\par
Lambertian radiance is view-independent. Projected area changes with viewpoint,\allowbreak{} and brightness still depends on the light–normal angle. Do not confuse radiance with the amount of energy collected from a fixed surface area.\par
{\small\itshape Topic: reflectance. Sources: HC2a pp.12–20.}

\end{minipage}
\par\vspace{12pt}
\begin{minipage}{\linewidth}
\textbf{Q05 — C}\par\smallskip
\textbf{C. The component of g perpendicular to that plane is unobservable.}\par
L has rank at most 2. Any vector in its nullspace can be added to g without changing Lg. More dependent observations cannot supply the missing component.\par
{\small\itshape Topic: photometric. Sources: HC2a pp.34–39; Lab1.}

\end{minipage}
\par\vspace{12pt}
\begin{minipage}{\linewidth}
\textbf{Q06 — D}\par\smallskip
\textbf{D. The normal is undefined; mask or otherwise handle the zero-reflectance/\allowbreak{}invalid pixel.}\par
A zero vector has no direction. Adding a small denominator can prevent a numerical exception,\allowbreak{} but does not recover a physically supported normal.\par
{\small\itshape Topic: photometric. Sources: Lab1; HC2a pp.35–45.}

\end{minipage}
\par\vspace{12pt}
\begin{minipage}{\linewidth}
\textbf{Q07 — C}\par\smallskip
\textbf{C. Interior intensities increase; 0 and 1 remain fixed.}\par
For r in (0,\allowbreak{}1),\allowbreak{} reducing a positive exponent below 1 makes \(r^{\gamma}\) larger than r. Always state the power-law convention; some software uses its reciprocal.\par
{\small\itshape Topic: histogram. Sources: HC2a pp.59–63; S1.}

\end{minipage}
\par\vspace{12pt}
\begin{minipage}{\linewidth}
\textbf{Q08 — B}\par\smallskip
\textbf{B. Different provisional labels can later be found to belong to one connected region.}\par
The first pass can encounter separate branches before their connecting pixels. The second pass resolves recorded label equivalences.\par
{\small\itshape Topic: morphology. Sources: HC2b pp.7–12.}

\end{minipage}
\par\vspace{12pt}
\begin{minipage}{\linewidth}
\textbf{Q09 — C}\par\smallskip
\textbf{C. Opening,\allowbreak{} then closing.}\par
Opening removes small foreground structures. Closing fills small background gaps/\allowbreak{}holes. Either can alter genuine features,\allowbreak{} so choose element size and shape with care.\par
{\small\itshape Topic: morphology. Sources: HC2b pp.13–45; S2.}

\end{minipage}
\par\vspace{12pt}
\begin{minipage}{\linewidth}
\textbf{Q10 — A}\par\smallskip
\textbf{A. Flip K horizontally and vertically before the sliding dot product.}\par
Convolution reverses both spatial offsets. Symmetric kernels conceal this distinction; asymmetric derivative or shift kernels do not.\par
{\small\itshape Topic: filtering. Sources: HC2b pp.78–80; S2.}

\end{minipage}
\par\vspace{12pt}
\begin{minipage}{\linewidth}
\textbf{Q11 — C}\par\smallskip
\textbf{C. 5}\par
Variances add: \(\sigma^2=3^2+4^2=25\). Standard deviations do not add directly.\par
{\small\itshape Topic: filtering. Sources: HC2b pp.82–87; S2.}

\end{minipage}
\par\vspace{12pt}
\begin{minipage}{\linewidth}
\textbf{Q12 — B}\par\smallskip
\textbf{B. \(\sqrt{a^2+b^2},\;\operatorname{arctan2}(b,a)\)}\par
Magnitude is distance from the origin. Phase is the quadrant-aware angle measured from the positive real axis; NumPy takes (imaginary,\allowbreak{} real). See SOURCE\_\allowbreak{}NOTES.md for the slide argument-order issue.\par
{\small\itshape Topic: fourier. Sources: HC2b pp.97–100; SOURCE\_\allowbreak{}NOTES.}

\end{minipage}
\par\vspace{12pt}
\begin{minipage}{\linewidth}
\textbf{Q13 — A}\par\smallskip
\textbf{A. Convolve the image with the corresponding derivative of the Gaussian.}\par
\(\frac{\partial(I*G)}{\partial x}=I*\frac{\partial G}{\partial x}\). Smoothing controls noise before measuring spatial changes.\par
{\small\itshape Topic: features. Sources: HC3a pp.20–54; S2; Lab2.}

\end{minipage}
\par\vspace{12pt}
\begin{minipage}{\linewidth}
\textbf{Q14 — D}\par\smallskip
\textbf{D. Keep the connected weak chain and discard the disconnected chain.}\par
Strong pixels seed edge following. Weak pixels survive only through a qualifying connection to a strong seed,\allowbreak{} under the declared neighbourhood rule.\par
{\small\itshape Topic: features. Sources: HC3a pp.58–65; S2.}

\end{minipage}
\par\vspace{12pt}
\begin{minipage}{\linewidth}
\textbf{Q15 — B}\par\smallskip
\textbf{B. Strong intensity variation in two independent directions: a good corner-like feature.}\par
Both displacement directions produce a large change. One large and one small eigenvalue instead indicates an edge; two small values indicate a flat patch.\par
{\small\itshape Topic: features. Sources: S3; Lab3 Sections 2–3.}

\end{minipage}
\par\vspace{12pt}
\begin{minipage}{\linewidth}
\textbf{Q16 — D}\par\smallskip
\textbf{D. A vertical line cannot be represented by a finite m; an orthogonal-distance line model avoids this problem.}\par
The slope-intercept model privileges y as the dependent variable. Total least squares uses perpendicular distances and a line representation that admits vertical lines.\par
{\small\itshape Topic: fitting. Sources: HC3a pp.76–79; S3.}

\end{minipage}
\par\vspace{12pt}
\begin{minipage}{\linewidth}
\textbf{Q17 — C}\par\smallskip
\textbf{C. \(I_t\) increases,\allowbreak{} often sharply.}\par
The all-inlier sample probability is \(w^s\). A lower w makes such samples rarer. RANSAC can work below 50\% inliers,\allowbreak{} but the required budget can be large.\par
{\small\itshape Topic: fitting. Sources: HC3a pp.91–109; S3.}

\end{minipage}
\par\vspace{12pt}
\begin{minipage}{\linewidth}
\textbf{Q18 — A}\par\smallskip
\textbf{A. \(A^{\mathsf T}A\): the local second-moment/\allowbreak{}structure tensor.}\par
\(A^{\mathsf T}A\) contains sums of \(I_x^2\),\allowbreak{} \(I_x I_y\) and \(I_y^2\). Its eigenvalues determine whether the local flow estimate is well constrained in both directions.\par
{\small\itshape Topic: flow. Sources: HC3b p.47; S3; Lab3.}

\end{minipage}
\par\vspace{12pt}
\begin{minipage}{\linewidth}
\textbf{Q19 — D}\par\smallskip
\textbf{D. A single common flow vector for all pixels in the window.}\par
A larger window adds equations but can mix incompatible motions. Local motion consistency,\allowbreak{} not simply pixel count,\allowbreak{} matters.\par
{\small\itshape Topic: flow. Sources: HC3b pp.39–47; Lab3.}

\end{minipage}
\par\vspace{12pt}
\begin{minipage}{\linewidth}
\textbf{Q20 — D}\par\smallskip
\textbf{D. Parallelism of lines.}\par
Affine transforms preserve parallelism and ratios along a line,\allowbreak{} not general angles or lengths. General homographies preserve straight lines but can map parallel lines to intersecting ones.\par
{\small\itshape Topic: geometry. Sources: HC4a pp.17–22; S4.}

\end{minipage}
\par\vspace{12pt}
\begin{minipage}{\linewidth}
\textbf{Q21 — C}\par\smallskip
\textbf{C. One homography,\allowbreak{} independent of scene depth.}\par
Pure rotation makes the image mapping depend on ray direction rather than depth: \(H=K_2 R_{21}K_1^{-1}\). A planar scene is the other important homography case.\par
{\small\itshape Topic: geometry. Sources: HC4a pp.33–35; S4.}

\end{minipage}
\par\vspace{12pt}
\begin{minipage}{\linewidth}
\textbf{Q22 — B}\par\smallskip
\textbf{B. Without a scale constraint,\allowbreak{} \(h=0\) is a trivial minimizer and does not define a homography.}\par
A homography is defined only up to nonzero scale. The constrained solution is the smallest-eigenvalue eigenvector of \(A^{\mathsf T}A\),\allowbreak{} equivalently the appropriate right singular vector of A.\par
{\small\itshape Topic: geometry. Sources: HC4a pp.51–55.}

\end{minipage}
\par\vspace{12pt}
\begin{minipage}{\linewidth}
\textbf{Q23 — B}\par\smallskip
\textbf{B. The match is ambiguous even if the nearest absolute distance is small.}\par
Repeated or weakly distinctive patterns can have several nearly equally plausible matches. The ratio compares distinctiveness relative to the next alternative.\par
{\small\itshape Topic: matching. Sources: HC4a pp.41–44; S4; Lab4.}

\end{minipage}
\par\vspace{12pt}
\begin{minipage}{\linewidth}
\textbf{Q24 — B}\par\smallskip
\textbf{B. Represent them identically,\allowbreak{} because it discards spatial arrangement.}\par
The vocabulary gives fixed-length count vectors. Plain BoVW is orderless; a spatial extension would be needed to retain layout.\par
{\small\itshape Topic: classification. Sources: HC4b pp.23–28; S4.}

\end{minipage}
\par\vspace{12pt}
\begin{minipage}{\linewidth}
\textbf{Q25 — D}\par\smallskip
\textbf{D. Fit the vocabulary on training descriptors; select settings on validation data; evaluate once on untouched test data.}\par
Unsupervised preprocessing can still leak information. The validation set selects hyperparameters; the test set measures the final chosen pipeline.\par
{\small\itshape Topic: classification. Sources: HC4b pp.25–28,\allowbreak{}55–58; Lab5.}

\end{minipage}
\par\vspace{12pt}
\begin{minipage}{\linewidth}
\textbf{Q26 — A}\par\smallskip
\textbf{A. 0}\par
\(I_t\) finds none of the minority examples: \(\mathrm{TP}=0\) and FN is the minority count. Accuracy alone can hide complete failure on a rare class.\par
{\small\itshape Topic: metrics. Sources: HC4b pp.80–86.}

\end{minipage}
\par\vspace{12pt}
\begin{minipage}{\linewidth}
\textbf{Q27 — D}\par\smallskip
\textbf{D. Translation equivariance: shifting the input shifts the feature map.}\par
Shared local operations follow translations. Invariance would mean unchanged output; pooling and aggregation can provide limited or additional invariance,\allowbreak{} not the same property.\par
{\small\itshape Topic: cnn. Sources: HC5b pp.22–24,\allowbreak{}64; S5.}

\end{minipage}
\par\vspace{12pt}
\begin{minipage}{\linewidth}
\textbf{Q28 — C}\par\smallskip
\textbf{C. Faster R-CNN uses a learned Region Proposal Network; Fast R-CNN normally relies on external proposals such as Selective Search.}\par
Fast R-CNN shares a full-image convolutional feature map but retains external proposals. Faster R-CNN also learns proposals through an RPN sharing backbone features.\par
{\small\itshape Topic: detection. Sources: HC5b pp.125–133; HC6a pp.18–38.}

\end{minipage}
\par\vspace{12pt}
\begin{minipage}{\linewidth}
\textbf{Q29 — A}\par\smallskip
\textbf{A. ResNet adds compatible activations; U-Net concatenates encoder and decoder features across matching scales.}\par
Residual addition supports identity-like mappings and gradient paths. U-Net concatenation transfers high-resolution spatial information to the decoder.\par
{\small\itshape Topic: architectures. Sources: HC6a pp.99–112,\allowbreak{}122–137.}

\end{minipage}
\par\vspace{12pt}
\begin{minipage}{\linewidth}
\textbf{Q30 — B}\par\smallskip
\textbf{B. An epipolar line on which the corresponding point must lie.}\par
Epipolar geometry reduces correspondence search from 2D to 1D. \(I_t\) does not select a unique point along the line. The nonzero condition excludes the epipole's degenerate zero output.\par
{\small\itshape Topic: stereo. Sources: HC6b pp.50–65.}

\end{minipage}
\par\vspace{12pt}
\begin{minipage}{\linewidth}
\textbf{Q31 — -270.00}\par\smallskip
\textbf{Submit -270.00.} Unrounded reference: \texttt{-270.0}.\par
\(X_w-C=\begin{bmatrix}2&3&4\end{bmatrix}^{\mathsf T}\),\allowbreak{} \(X_c=\begin{bmatrix}-3&2&4\end{bmatrix}^{\mathsf T}\). \(q_0=800(-3)+20\cdot2+320\cdot4=-1080\). Divide by 4: \(u=-270\). A negative pixel coordinate can simply be outside the field of view.\par
{\small\itshape Topic: camera. Sources: HC1b pp.24–45; S1.}

\end{minipage}
\par\vspace{12pt}
\begin{minipage}{\linewidth}
\textbf{Q32 — 0.75}\par\smallskip
\textbf{Submit 0.75.} Unrounded reference: \texttt{0.7452154112317649}.\par
\(L^{\mathsf T}L=\operatorname{diag}(1.28,1.28,2.44)\) and \(L^{\mathsf T}i=\begin{bmatrix}0.384&0.064&1.66\end{bmatrix}^{\mathsf T}\). Thus \(g=\begin{bmatrix}0.3&0.05&1.66/2.44\end{bmatrix}^{\mathsf T}\). Its norm is about 0.7452,\allowbreak{} so submit 0.75. Do not normalize g before taking its magnitude.\par
{\small\itshape Topic: photometric. Sources: HC2a pp.34–39; Lab1; Coding examples Q1,\allowbreak{}Q5.}

\end{minipage}
\par\vspace{12pt}
\begin{minipage}{\linewidth}
\textbf{Q33 — 4.00}\par\smallskip
\textbf{Submit 4.00.} Unrounded reference: \texttt{4.0}.\par
\(p=-\frac{2}{2}=-1\) and \(q=-\frac{-1}{2}=0.5\). \(z(3,4)=5+3(-1)+4\cdot0.5=4\). The declared sign convention is essential.\par
{\small\itshape Topic: photometric. Sources: HC2a pp.41–43; Lab1; SOURCE\_\allowbreak{}NOTES.}

\end{minipage}
\par\vspace{12pt}
\begin{minipage}{\linewidth}
\textbf{Q34 — 1}\par\smallskip
\textbf{Submit 1.} Unrounded reference: \texttt{1}.\par
\(\operatorname{CDF}(1)=\frac{1+3}{8}=0.5\); \(\lfloor3\cdot0.5\rfloor=1\). Other equalization conventions can differ,\allowbreak{} which is why the mapping is specified.\par
{\small\itshape Topic: histogram. Sources: HC2a pp.51–63; S1.}

\end{minipage}
\par\vspace{12pt}
\begin{minipage}{\linewidth}
\textbf{Q35 — 8}\par\smallskip
\textbf{Submit 8.} Unrounded reference: \texttt{8}.\par
Dilation gives \(\begin{bmatrix}1&1&1&1&1&1&1&1&1&1&0\end{bmatrix}\). Erosion of that fixed-length result gives \(\begin{bmatrix}0&1&1&1&1&1&1&1&1&0&0\end{bmatrix}\),\allowbreak{} which has 8 foreground entries.\par
{\small\itshape Topic: morphology. Sources: HC2b pp.13–45; S2.}

\end{minipage}
\par\vspace{12pt}
\begin{minipage}{\linewidth}
\textbf{Q36 — 8}\par\smallskip
\textbf{Submit 8.} Unrounded reference: \texttt{8}.\par
The flipped kernel is
\[\begin{bmatrix}-2&0&3\\0&-1&0\\2&0&1\end{bmatrix}.\]
 The dot product is \(-2+1+4+5=8\). The unflipped result would be different.\par
{\small\itshape Topic: filtering. Sources: HC2b pp.78–80; S2; Lab2.}

\end{minipage}
\par\vspace{12pt}
\begin{minipage}{\linewidth}
\textbf{Q37 — 2.83}\par\smallskip
\textbf{Submit 2.83.} Unrounded reference: \texttt{2.8284271247461903}.\par
\(F_1=1-2\mathrm{i}-3+4\mathrm{i}=-2+2\mathrm{i}\). Its magnitude is \(\sqrt{4+4}=\sqrt{8}\),\allowbreak{} so submit 2.83.\par
{\small\itshape Topic: fourier. Sources: HC2b pp.95–104.}

\end{minipage}
\par\vspace{12pt}
\begin{minipage}{\linewidth}
\textbf{Q38 — 59.04}\par\smallskip
\textbf{Submit 59.04.} Unrounded reference: \texttt{59.04}.\par
\(\det(M)=81-9=72\); \(\operatorname{tr}(M)=18\). \(R=72-0.04\cdot324=59.04\). A positive response alone still needs thresholding and spatial non-maximum suppression.\par
{\small\itshape Topic: features. Sources: S3; Lab3.}

\end{minipage}
\par\vspace{12pt}
\begin{minipage}{\linewidth}
\textbf{Q39 — 2.36}\par\smallskip
\textbf{Submit 2.36.} Unrounded reference: \texttt{2.3636363636363638}.\par
\(A^{\mathsf T}A\) =\allowbreak{}
\[\begin{bmatrix}6&1\\1&2\end{bmatrix},\]
 \(A^{\mathsf T}(-I_t)=\begin{bmatrix}13&0\end{bmatrix}^{\mathsf T}\). Solve \(6u+v=13\) and \(u+2v=0\): \(u=\frac{26}{11}\),\allowbreak{} \(v=-\frac{13}{11}\). Submit 2.36; forgetting the minus sign reverses the estimate.\par
{\small\itshape Topic: flow. Sources: HC3b pp.39–47; S3; Lab3.}

\end{minipage}
\par\vspace{12pt}
\begin{minipage}{\linewidth}
\textbf{Q40 — 34}\par\smallskip
\textbf{Submit 34.} Unrounded reference: \texttt{34}.\par
The all-inlier probability is \(0.6^4=0.1296\). The ratio of logarithms is approximately 33.18; round upwards,\allowbreak{} not to the nearest integer. \(N=34\).\par
{\small\itshape Topic: fitting. Sources: HC3a pp.104–105; HC4a pp.45–48.}

\end{minipage}
\par\vspace{12pt}
\begin{minipage}{\linewidth}
\textbf{Q41 — 7.33}\par\smallskip
\textbf{Submit 7.33.} Unrounded reference: \texttt{7.333333333333333}.\par
\(H\begin{bmatrix}5&4&1\end{bmatrix}^{\mathsf T}=\begin{bmatrix}11&2&1.5\end{bmatrix}^{\mathsf T}\). Dehomogenize: \(x'=\frac{11}{1.5}=7.3333\),\allowbreak{} so submit 7.33. The third coordinate is not generally 1.\par
{\small\itshape Topic: geometry. Sources: HC4a pp.20–22,\allowbreak{}51; S4.}

\end{minipage}
\par\vspace{12pt}
\begin{minipage}{\linewidth}
\textbf{Q42 — 48.75}\par\smallskip
\textbf{Submit 48.75.} Unrounded reference: \texttt{48.75}.\par
Weights are 0.1875,\allowbreak{} 0.0625,\allowbreak{} 0.5625 and 0.1875. The weighted sum is \(1.875+1.875+28.125+16.875=48.75\). Do not interchange x and y.\par
{\small\itshape Topic: geometry. Sources: HC4a pp.24–31; Lab4.}

\end{minipage}
\par\vspace{12pt}
\begin{minipage}{\linewidth}
\textbf{Q43 — 0.63}\par\smallskip
\textbf{Submit 0.63.} Unrounded reference: \texttt{0.631578947368421}.\par
\(\mathrm{TP}=12,\;\mathrm{FP}=2+4=6,\;\mathrm{FN}=3+5=8\). \(F_1=\frac{2\mathrm{TP}}{2\mathrm{TP}+\mathrm{FP}+\mathrm{FN}}=\frac{24}{38}=0.6315789\). Submit 0.63.\par
{\small\itshape Topic: metrics. Sources: HC4b pp.80–86; S4.}

\end{minipage}
\par\vspace{12pt}
\begin{minipage}{\linewidth}
\textbf{Q44 — 0.57}\par\smallskip
\textbf{Submit 0.57.} Unrounded reference: \texttt{0.5666666666666667}.\par
\(\mathrm{AP}=\frac{1/1+2/3+3/5}{4}=\frac{17}{30}\). The missing relevant image contributes zero,\allowbreak{} so the denominator is 4,\allowbreak{} not 3. Submit 0.57.\par
{\small\itshape Topic: metrics. Sources: HC5a pp.19–24,\allowbreak{}68–78; S4.}

\end{minipage}
\par\vspace{12pt}
\begin{minipage}{\linewidth}
\textbf{Q45 — 0.33}\par\smallskip
\textbf{Submit 0.33.} Unrounded reference: \texttt{0.3333333333333333}.\par
Each area is 24. The intersection is width 4 times height 3 =\allowbreak{} 12. \(\text{Union}=24+24-12=36\). \(\mathrm{IoU}=\frac{12}{36}=\frac{1}{3}\); submit 0.33.\par
{\small\itshape Topic: detection. Sources: HC5a p.62; S4.}

\end{minipage}
\par\vspace{12pt}
\begin{minipage}{\linewidth}
\textbf{Q46 — 1680}\par\smallskip
\textbf{Submit 1680.} Unrounded reference: \texttt{1680}.\par
\(H_{\mathrm{out}}=\left\lfloor\frac{31+2-5}{2}\right\rfloor+1=15\); \(W_{\mathrm{out}}=\left\lfloor\frac{29-3}{2}\right\rfloor+1=14\). \(\text{Total}=15\cdot14\cdot8=1680\). Count output positions,\allowbreak{} not input positions.\par
{\small\itshape Topic: cnn. Sources: HC5b pp.25–43,\allowbreak{}66–68; S5.}

\end{minipage}
\par\vspace{12pt}
\begin{minipage}{\linewidth}
\textbf{Q47 — 368}\par\smallskip
\textbf{Submit 368.} Unrounded reference: \texttt{368}.\par
Each output channel has \(5\cdot3\cdot3=45\) weights and one bias. There are 8 output channels,\allowbreak{} so \(8(45+1)=368\). Spatial resolution and stride do not change this count.\par
{\small\itshape Topic: cnn. Sources: HC5b pp.66–73; S5.}

\end{minipage}
\par\vspace{12pt}
\begin{minipage}{\linewidth}
\textbf{Q48 — 7850}\par\smallskip
\textbf{Submit 7850.} Unrounded reference: \texttt{7850}.\par
The flattened input has \(7\cdot7\cdot16=784\) values. Count \(784\cdot10\) weights plus \(10\text{ biases}=7850\).\par
{\small\itshape Topic: cnn. Sources: HC5b pp.16–24; S5.}

\end{minipage}
\par\vspace{12pt}
\begin{minipage}{\linewidth}
\textbf{Q49 — 8}\par\smallskip
\textbf{Submit 8.} Unrounded reference: \texttt{8}.\par
After conv: \(r=3,\;j=1\). After pool: \(r=4,\;j=2\). After second conv: \(r=4+2\cdot2=8,\;j=2\). Pooling changes the spacing used by later layers.\par
{\small\itshape Topic: cnn. Sources: HC5b pp.80–85; HC6a pp.6–9.}

\end{minipage}
\par\vspace{12pt}
\begin{minipage}{\linewidth}
\textbf{Q50 — 2.00}\par\smallskip
\textbf{Submit 2.00.} Unrounded reference: \texttt{2.0}.\par
\(l_2=\begin{bmatrix}0&-1&7\end{bmatrix}^{\mathsf T}\),\allowbreak{} so the line is \(y=7\). The candidate has \(y=9\),\allowbreak{} giving \(\text{distance}\;\frac{\lvert-9+7\rvert}{1}=2\). This tests the line constraint,\allowbreak{} not whether the candidate is the true appearance match.\par
{\small\itshape Topic: stereo. Sources: HC6b pp.50–65.}

\end{minipage}
\par\vspace{12pt}

\end{document}
